\documentclass[12pt]{article}
\usepackage[margin=1in]{geometry}
\usepackage{amsmath,amssymb}
\usepackage{natbib}
\usepackage{enumitem}

\title{Referee Report: ``Quantity-Space Prior Specification for Bayesian
Hierarchical Models: Theory and Applications to Risk Preferences''}
\author{Simulated Referee Report\\
\small Prepared for internal review, March 2026}
\date{}

\begin{document}
\maketitle

\section*{Summary}

This paper develops a framework for specifying hyperpriors in Bayesian
hierarchical models (BHMs) through two contributions: (i) quantity-space prior
specification, where priors are stated on predictive quantities (certainty
equivalents, choice probabilities) rather than structural parameters, using a
Jacobian correction; and (ii) a distribution-free variance decomposition that
reduces the remaining degrees of freedom to two interpretable parameters,
$\rho$ (variance partition ratio) and $\kappa$ (coefficient of variation of
heterogeneity). The framework is demonstrated with rank-dependent utility (RDU)
estimation on lottery choice data from Harrison \& Ng (2016).

The paper addresses a genuine practical problem in applied Bayesian work and
offers a clean analytical solution. The theoretical development is rigorous,
the application is informative, and the empirical findings---particularly the
lottery-specific risk attitude patterns and the prior sensitivity
analysis---are a valuable contribution in their own right.

\section*{Major Comments}

\subsection*{1. The comparison is not fully symmetric}

The paper compares an informative CE-space model (Beta-Beta hierarchy with
Jacobian correction) against a diffuse parameter-space model (Normal--Inverse-Gamma
with Gao et al.\ 2022 hyperpriors). These are not two specifications of the
\textit{same} model with different priors; they are structurally different
models with different likelihood contributions (the Jacobian term is present
in one and absent in the other). A cleaner comparison would hold the model
structure fixed and vary only the hyperprior. For example, the author could
estimate the CE-space Beta-Beta model with both informative and diffuse
Beta-Beta hyperpriors, or estimate the NIG model with both informative and
diffuse NIG hyperpriors. The current comparison confounds the effect of the
prior specification method with the effect of the distributional family.

\subsection*{2. Independence of quantities in the direct prior}

The framework specifies independent marginal priors over the $p$ predictive
quantities. For the RDU application, this means $\text{CE}_1$, $\text{CE}_2$,
$\text{CE}_3$, and $P$ are treated as independent in the prior. But these
quantities are deterministic functions of the same parameters $(r, \phi, \eta,
\lambda)$, so they are necessarily dependent. The paper should discuss what
constraints the independent-marginals assumption places on the choice of
reference lotteries. Intuitively, if two reference lotteries are very similar,
their CEs will be nearly perfectly correlated, and specifying independent
priors on them could be highly informative in ways the researcher does not
intend. The no-FOSD condition is necessary for this reason, but is it
sufficient?

\subsection*{3. Moment matching is only approximate for the Beta-Beta hierarchy}

The variance decomposition (equations 7--8) is exact for any hierarchy where
$\psi_m$ bijects to the conditional mean and $\psi_v$ bijects to the
conditional variance. In the Beta-Beta model used in the application, the
conditional variance is $\mu_k(1 - \mu_k)\psi_k$, which depends on both
$\mu_k$ and $\psi_k$. The decomposition in equation (25) uses
$\mathbb{E}[\mu_k(1-\mu_k)\psi_k] = (1-\rho_k)v^*_k$, which requires
independence of $\mu_k$ and $\psi_k$ to factor the expectation. This
independence is assumed by the Beta-Beta specification (separate hyperpriors on
$\mu_k$ and $\psi_k$) but is not a consequence of the general framework.
The paper should be more explicit that the general variance decomposition
requires a choice of parameterization where the mean and variance
hyperparameters enter independently, and that the $\psi_k$ parameterization in
the Beta-Beta model is specifically chosen to achieve this.

\subsection*{4. Non-centered parameterization is not available}

The paper uses centered parameterization for the CE-space model, storing
individual parameters $(r, \phi, \eta, \lambda)$ directly in the chain state.
The population prior is evaluated in CE space via the forward map and Jacobian
correction. This means the standard non-centered parameterization (NCP) trick
is not available, because it would require inverting the CE mapping at every
MCMC step. The paper does not discuss this computational trade-off. For the NIG
comparison model, NCP \textit{is} available (and was used in previous versions
of this work). The centered parameterization is known to create funnel
geometries that impede mixing when population variance is small. Given that the
paper reports convergence issues for $\sim$18\% of subjects, this trade-off
deserves discussion. Are the convergence difficulties concentrated in subjects
where $\psi_k$ is small?

\subsection*{5. Identification and the 17\% finding}

The paper reports that approximately 17--18\% of subjects have opposing
posterior inferences under the two models, with the population prior being the
deciding factor. The posterior decomposition for Subject 74 is informative:
the likelihood ratio is 29:1 in favor of the risk-seeking mode, but the
population prior is $10^4$:1 against it. This is a strong prior.

The paper presents this as ``precisely the role an informative prior should
play.'' While this may be true for subjects with genuinely uninformative data,
the finding also raises questions about whether the informative prior is
\textit{too strong} for these subjects. A sensitivity analysis varying $\rho$
and $\kappa$ and examining how the fraction of prior-dominated subjects changes
would be informative. If modest changes in $\kappa$ move the 17\% substantially,
this would suggest the results are fragile.

\section*{Minor Comments}

\begin{enumerate}[leftmargin=*]

\item \textbf{Section 2}: The literature review of belief elicitation
(paragraphs on PSRs, Hossain \& Graves, Schlag, etc.) is thorough but
tangential. The paper is about \textit{researcher} prior specification, not
\textit{subject} belief elicitation. The connection between the two is noted
but could be made more concisely. This section could be shortened by half
without losing the paper's contribution.

\item \textbf{Section 3.2, equation (6)}: The notation $\psi_m$ and $\psi_v$
for the mean and variance hyperparameters is potentially confusing given that
$\psi_k$ is later used for the Beta-Beta dispersion parameter. Consider
different notation for the general hyperparameters.

\item \textbf{Section 3.3}: The NIG example and sensitivity heatmaps are
well-done. One suggestion: the heatmap figure caption could note that these
are computed analytically (not via MCMC), since this is a strength of the
framework---the sensitivity analysis requires no estimation.

\item \textbf{Section 4}: The paper uses ``observable quantities'' in some
places and ``predictive quantities'' in others. The terminology should be
consistent throughout. (``Predictive quantities'' appears to be the intended
term.)

\item \textbf{Section 5, equation (15)}: The comparison prior
$\mu^j \sim N(0, 100)$, $\sigma^{2,j} \sim \text{IG}(0.001, 0.001)$ is
presented without discussion of what this implies for predictive quantities.
Given the paper's emphasis on quantity-space reasoning, computing the prior
predictive CEs implied by the diffuse prior would strengthen the comparison.
What does the Gao et al.\ prior predict about CEs? If it puts substantial
mass on extreme or nonsensical values, this would reinforce the paper's
argument.

\item \textbf{Section 5, Results}: The footnote on MCMC details is sparse.
For a paper that relies heavily on MCMC results, more detail on convergence
diagnostics (R-hat distribution, effective sample sizes, acceptance rates)
is warranted, at minimum in an appendix.

\item \textbf{Section 5, Figure 3 (CE overlay)}: The ``hyperprior predictive''
curve for the informative model is drawn from the Beta-Beta hyperprior, while
the ``diffuse posterior'' is drawn by computing CEs from NIG hyper-posterior
draws via the forward map. These use different computational procedures. The
paper should note whether the two posterior predictive curves are computed
comparably (both via forward map? both via direct hierarchy draws?).

\item \textbf{Section 5, CDF difference figure}: This is an effective
visualization. Consider adding a panel or inset showing the distribution of
the average absolute median CE difference across subjects, so the reader can
see where the thresholds (2, 10) fall relative to the full distribution.

\item \textbf{Section 6, Limitations}: The fourth limitation (identification
with 80 choices) could be strengthened. Is the fraction of prior-dominated
subjects likely to decrease with more observations per subject? A brief power
analysis or simulation would be informative, though not strictly necessary for
this paper.

\item \textbf{Conclusion}: The paper claims the framework ``reduces'' the
researcher's burden. The revised wording (``replaces opaque with
interpretable'') is more accurate, but the claim could be further sharpened:
the key advantage is not just interpretability but
\textit{consistency-by-construction} with the prior predictive distribution,
which direct hyperprior specification does not guarantee.

\item \textbf{Appendix A (Full Covariance)}: The relevance of this appendix
to the current paper is unclear. The paper uses independent Beta-Beta
hierarchies, not the NIW model derived in Appendix A. If this appendix is
intended for a future extension, it should be noted; otherwise it could be
removed.

\item \textbf{Missing}: The paper does not report posterior summaries for the
structural parameters $(r, \phi, \eta, \lambda)$. While the emphasis on
quantity-space is the point, many readers will want to know the estimated
population mean and heterogeneity of risk aversion, probability weighting,
etc. A posterior summary table would be a valuable addition.

\item \textbf{Missing}: There is no discussion of model fit. Do subjects'
observed choice patterns match the model's predictions? A posterior predictive
check on choice probabilities (not CEs) would address whether the RDU model
with the estimated parameters actually explains the data well.

\end{enumerate}

\section*{Overall Assessment}

The paper makes two clean, well-motivated methodological contributions and
demonstrates them with an informative application. The quantity-space prior
specification and the $(\rho, \kappa)$ variance decomposition are useful tools
that should see adoption beyond the risk preference setting. The empirical
findings, particularly the lottery-specific risk attitude patterns and the
distributional prior sensitivity analysis, add genuine value.

The main weaknesses are the asymmetric comparison (mixing distributional family
with prior specification), the limited discussion of computational trade-offs
(NCP unavailability, centered parameterization), and the absence of posterior
summaries for the structural parameters. These are addressable in revision.

I recommend revision and resubmission.

\end{document}
